\section{Implantação}
\label{sec:implantacao}

A implantação de um sistema de software em um ambiente de produção é a etapa final que torna a aplicação acessível e utilizável. Para um sistema distribuído como a plataforma de processamento de mídias, a estratégia de implantação deve garantir não apenas a disponibilidade, mas também a escalabilidade, a segurança e a consistência do ambiente.

Esta seção detalha o plano de implantação do projeto, que é baseado em uma abordagem moderna de containerização com Docker e orquestração com Docker Compose. A estratégia adotada visa criar um ambiente de produção portátil, replicável e otimizado, separando claramente as configurações de desenvolvimento das de produção para garantir estabilidade e performance. A subseção a seguir descreve a arquitetura de implantação, incluindo a configuração dos contêineres, a rede e o uso de servidores de aplicação robustos.
